%%%%%%%%%%%%%%%%%%%%%%%%%%%%%%%%%%%%%%%%%%%%%%%%%%%%%%%%%%%%%%%%%%%%%%%%%%%%%%%%

\pagestyle{empty}
\begin{abstract}
				En ámbitos críticos como a imaxe médica non abonda con que un modelo acerte: o profesional precisa entender por que decide así, e a normativa europea esixe transparencia nos sistemas de alto risco. As explicacións \emph{a posteriori} adoitan ser parciais, polo que unha arquitectura interpretable e explicable por construción ofrece unha vía prometedora. \textsc{Crate} é un transformer de visión no que cada compoñente deriva do obxectivo \acrfull{mcr2} disperso, e este traballo estuda como ampliar esas propiedades e levalas á análise de imaxes de retina. Proponse unha variante de segunda orde da atención, o cruce de arquitecturas e a compartición dos tensores $U$, de proxección de subespazos e $D$, dicionario de patróns entre capas, e estendense as técnicas de explicabilidade con mapas de atención de segunda orde e de maior resolución para estudar $U$ e coa recuperación de parches representativos para estudar $D$. A experimentación en conxuntos de imaxes xerais, como ImageNet e de imaxes de fondo de ollo, como DRIVE e RFMiD, inclúe ata onde sabemos, a primeira aplicación de \textsc{Crate} a imaxe de fondo de retina. Os resultados son alentadores: compartir $U$ e $D$ mantén un rendemento comparable cunha rede máis pequena e, en RFMiD, produce mapas de atención máis variados e informativos; a variante de segunda orde, aínda sen superar en xeral ao modelo orixinal, revela o rol das cabezas, e o fracaso do cruce confirma devandita interpretación. As propiedades de \textsc{Crate} mantéñense así no dominio da retina, como paso previo á súa validación con profesionais clínicos.
  \vspace*{25pt}
  \begin{segundoresumo}
					In critical domains such as medical imaging, accuracy is not enough: clinicians need to understand why a model decides as it does, and European regulation demands transparency in high-risk systems. \emph{Post-hoc} explanations are often partial, so an architecture that is interpretable and explainable by construction offers a promising alternative. \textsc{Crate} is a vision transformer in which each component derives from the sparse \acrfull{mcr2} objective, and this work studies how to extend its properties and bring them to the analysis of retinal images. We propose a second-order attention variant, architecture crossing, and sharing the $U$ (subspace projection) and $D$ (pattern dictionary) tensors across layers, and extend explainability techniques with second-order and higher-resolution attention maps to study $U$ and with the retrieval of representative patches to study $D$. Experiments on general image datasets such as ImageNet and on ocular fundus images such as DRIVE and RFMiD include, to our knowledge, the first application of \textsc{Crate} to retinal fundus images. Results are encouraging: sharing $U$ and $D$ keeps comparable performance with a smaller network and, on RFMiD, yields more varied and informative attention maps; the second-order variant, while not generally outperforming the original, reveals the role of the heads, and the failure of crossing supports that reading. \textsc{Crate}'s properties thus carry over to the retinal domain, a first step before validation with clinicians.
  \end{segundoresumo}
\vspace*{25pt}
\begin{multicols}{2}
\begin{description}
\item [\palabraschaveprincipal:] \mbox{} \\[-20pt]
\begin{itemize}
  \item \textsc{Crate}
  \item Interpretabilidade
  \item Explicabilidade
  \item Aprendizaxe desenrolada
  \item MCR$^2$ disperso
  \item Imaxe médica
  \item Fondo de retina
\end{itemize}
\end{description}
\begin{description}
\item [\palabraschavesecundaria:] \mbox{} \\[-20pt]
\begin{itemize}
  \item \textsc{Crate}
  \item Interpretability
  \item Explainability
  \item Algorithm unrolling
  \item Sparse MCR$^2$
  \item Medical imaging
  \item Retinal fundus
\end{itemize}
\end{description}
\end{multicols}

\end{abstract}
\pagestyle{fancy}

%%%%%%%%%%%%%%%%%%%%%%%%%%%%%%%%%%%%%%%%%%%%%%%%%%%%%%%%%%%%%%%%%%%%%%%%%%%%%%%%
